\chapter{广义 Verma 模}
\label{ch:06.6}

\section{Weyl 群中的陪集}
\label{sec:06.6.1}

设 $\Sigma$ 是约化根系，$\Sigma^+$ 是一组正根。记 $\Pi$ 为相应的单根集。照旧置
$$\Sigma^+_w=\Sigma^+\cap\{-w^{-1}(\Sigma^+)\}=\{\alpha\in\Sigma^+\mid w\alpha\in-\Sigma^+\}.$$
\begin{annote}
本章改用抽象根系的语言。这里的 $\Sigma$ 就是第 5 章 §1 中的那个根系；在那一章，经 Borel 子代数 $\bb_0$ 决定的特殊化，它在 $\h_0^*$ 中的像记作 $R_0$。本章不再特化，直接写 $\Sigma$，并把它的 Weyl 群记作 $W$。两者是同一根系在不同语境下的记号。[来源：底本 p.113 与 p.93，两处开头的对照]
\end{annote}

\begin{lemma}\label{lem:06.1.1}
设 $w\in W$，$\alpha\in\Pi$。则下列条件等价：
\begin{enumerate}[label=(\roman*)]
\item $\ell(ws_\alpha)=\ell(w)+1$；
\item $\alpha\notin\Sigma^+_w$。
\end{enumerate}
\end{lemma}

\begin{proof}
设 $\ell(ws_\alpha)=\ell(w)+1$。则由 L.3.12.(ii) 有
$$\Sigma^+_{ws_\alpha}=s_\alpha(\Sigma^+_w)\cup\{\alpha\},$$
即 $\alpha\in\Sigma^+_{ws_\alpha}$。这说明
$$-w\alpha=ws_\alpha\alpha\in-\Sigma^+,$$
即 $\alpha\notin\Sigma^+_w$。

若 $\ell(ws_\alpha)=\ell(w)-1$，则对 $w'=ws_\alpha$ 有 $\ell(w's_\alpha)=\ell(w')+1$，且 $\alpha\notin\Sigma^+_{w'}$。这又说明
$$w\alpha=w's_\alpha\alpha=-w'\alpha\in-\Sigma^+,$$
从而 $\alpha\in\Sigma^+_w$。
\end{proof}

设 $\Theta\subset\Pi$。记 $\Sigma_\Theta$ 为 $\Sigma$ 中由 $\Theta$ 生成的子根系，$W_\Theta$ 为 $W$ 中由单反射 $S_\Theta=\{s_\alpha\mid\alpha\in\Theta\}$ 生成的子群。显然，$(W_\Theta,S_\Theta)$ 的长度函数就是 $\ell$ 在 $W_\Theta$ 上的限制。再置
$$W^\Theta=\{w\in W\mid\Sigma^+_w\cap\Theta=\emptyset\}=\{w\in W\mid\Theta\subset w^{-1}(\Sigma^+)\}.$$

\begin{theorem}\label{thm:06.1.2}
每个 $w\in W$ 都有唯一分解 $w=w't$，其中 $w'\in W^\Theta$，$t\in W_\Theta$。此外 $\ell(w)=\ell(w')+\ell(t)$。
\end{theorem}

\begin{proof}
由 1.1，对任意 $w\in W$，下列条件等价：
\begin{enumerate}[label=(\roman*)]
\item $w\in W^\Theta$；
\item 对任意 $\alpha\in\Theta$ 有 $\ell(ws_\alpha)=\ell(w)+1$。
\end{enumerate}

先断言：左 $W_\Theta$-陪集中的最短元必属于 $W^\Theta$。设 $w$ 是某个左 $W_\Theta$-陪集中的最短元，且 $w\notin W^\Theta$。则存在 $\alpha\in\Theta$ 使 $\ell(ws_\alpha)=\ell(w)-1$。于是同一个左陪集中有长度更短的元，与假设矛盾。

下面对长度作归纳，证明左 $W_\Theta$-陪集中的每个元 $w$ 都有形如 $w=w't$ 的分解，其中 $w'\in W^\Theta$、$t\in W_\Theta$，且 $\ell(w)=\ell(w')+\ell(t)$。最短元的情形已经证过。取左陪集中任一 $w$。若 $w\in W^\Theta$，则已了。若 $w\notin W^\Theta$，由上述注记，可取 $\alpha\in\Theta$ 使 $\ell(ws_\alpha)=\ell(w)-1$。由归纳假设，$ws_\alpha=w't'$，其中 $w'\in W^\Theta$、$t'\in W_\Theta$，且 $\ell(ws_\alpha)=\ell(w')+\ell(t')$。于是 $w=w't's_\alpha$。置 $t=t's_\alpha$。则 $w=w't$，$t\in W_\Theta$，且 $\ell(t)\le\ell(t')+1$。又有
$$\ell(w)\le\ell(w')+\ell(t)\le\ell(w')+\ell(t')+1=\ell(ws_\alpha)+1=\ell(w),$$
故 $\ell(w)=\ell(w')+\ell(t)$。这就证明了分解的存在性。

为证唯一性，只需说明每个左 $W_\Theta$-陪集中至多有一个 $W^\Theta$ 的元。设 $w,w'\in W^\Theta$，且 $w=w't$，$t\in W_\Theta$。则
$$\Theta\subset w^{-1}(\Sigma^+)=t^{-1}w'^{-1}(\Sigma^+).$$
集 $\Sigma^+_\Theta=\Sigma_\Theta\cap\Sigma^+$ 是 $\Sigma_\Theta$ 中由 $\Theta$ 决定的正根集。于是
$$\Sigma^+_\Theta\subset t^{-1}w'^{-1}(\Sigma^+),$$
从而
$$t(\Sigma^+_\Theta)\subset w'^{-1}(\Sigma^+).$$
同理
$$\Sigma^+_\Theta\subset w'^{-1}(\Sigma^+).$$
因 $W_\Theta$ 同构于 $\Sigma_\Theta$ 的 Weyl 群，若 $t\neq 1$，就存在根 $\beta\in\Sigma^+_\Theta$ 使 $-\beta\in t(\Sigma^+_\Theta)$。这会推出 $\beta,-\beta\in w'^{-1}(\Sigma^+)$，不可能。故 $t=1$，$w=w'$。
\end{proof}

设 $w_\Theta$ 是 $W_\Theta$ 中的最长元。这个元由下述性质刻画。

\begin{lemma}\label{lem:06.1.3}
$w_\Theta$ 是 $W$ 中唯一具有下述性质的元：
\begin{enumerate}[label=(\roman*)]
\item $w_\Theta(\Theta)=-\Theta$；
\item $w_\Theta$ 置换 $\Sigma^+_\Theta$ 之外的正根。
\end{enumerate}
\end{lemma}

\begin{proof}
为证 $w_\Theta$ 满足 (i)，只需注意它把 $\Sigma_\Theta$ 中的正根映为负根。

设 $\beta\in\Theta$。则反射 $s_\beta$ 置换 $\Sigma^+-\{\beta\}$ 中的元，也置换 $\Sigma_\Theta$ 中的根。这说明它也置换 $\Sigma^+_\Theta$ 之外的正根。对长度作归纳，即知 $W_\Theta$ 的每个元都置换 $\Sigma^+_\Theta$ 之外的正根。故 $w_\Theta$ 满足 (ii)。

反之，若 $w$ 满足 (i) 与 (ii)，则
$$w(\Sigma^+)=(-\Sigma^+_\Theta)\cup(\Sigma^+-\Sigma^+_\Theta).$$
由于 $W$ 在 $\Sigma$ 的全部正根系构成的集上单可迁地作用，这样的 $w\in W$ 只有一个。
\end{proof}

\begin{theorem}\label{thm:06.1.4}
\begin{enumerate}[label=(\roman*)]
\item $W$ 中每个左 $W_\Theta$-陪集有唯一的最短元，它属于 $W^\Theta$。
\item 若 $w$ 是左 $W_\Theta$-陪集 $C$ 的最短元，则 $ww_\Theta$ 是该陪集中唯一的最长元。
\item $W$ 中每个右 $W_\Theta$-陪集有唯一的最短元。
\item 若 $w$ 是右 $W_\Theta$-陪集 $C$ 的最短元，则 $w_\Theta w$ 是该陪集中唯一的最长元。
\end{enumerate}
\end{theorem}

\begin{proof}
(i) 与 (ii) 由 1.2 立见。$W$ 的反自同构 $w\mapsto w^{-1}$ 保持 $W_\Theta$、$w_\Theta$ 与长度函数 $\ell:W\to\Z_+$，并把左 $W_\Theta$-陪集映为右 $W_\Theta$-陪集。故由 (i)、(ii) 即得 (iii)、(iv)。
\end{proof}

因此，$W^\Theta$ 是左 $W_\Theta$-陪集的一个截面，由各陪集的最短元组成。于是，右 $W_\Theta$-陪集的最短元构成集
$$\{w\in W\mid w^{-1}\in W^\Theta\}=\{w\in W\mid\Theta\subset w(\Sigma^+)\}.$$
这给出下述结果。

\begin{lemma}\label{lem:06.1.5}
集
$${}^{\Theta}W=\{w\in W\mid\Theta\subset-w(\Sigma^+)\}$$
是右 $W_\Theta$-陪集的一个截面，由各陪集的最长元组成。
\end{lemma}

\begin{proof}
设 $w$ 是某右 $W_\Theta$-陪集的最短元。则由上面的讨论，$\Theta\subset w(\Sigma^+)$。于是
$$-\Theta=w_\Theta(\Theta)\subset w_\Theta w(\Sigma^+),$$
故该右陪集的最长元 $w_\Theta w$ 属于 ${}^{\Theta}W$。

反之，若 $w\in{}^{\Theta}W$，则 $\Theta\subset-w(\Sigma^+)$，且
$$-\Theta=w_\Theta(\Theta)\subset-w_\Theta w(\Sigma^+),$$
这说明 $w_\Theta w$ 是其右 $W_\Theta$-陪集中的最短元。故 $w$ 是该陪集的最长元。
\end{proof}

对 $W$ 的右 $W_\Theta$-陪集 $C$，记 $w_C$ 为 ${}^{\Theta}W$ 中相应的元。把 $W$ 上 Bruhat 序在 ${}^{\Theta}W$ 上诱导的序关系搬过来，就在全体右 $W_\Theta$-陪集构成的集 $W_\Theta\backslash W$ 上定义了一个序关系。

\begin{proposition}\label{prop:06.1.6}
设 $C$ 是 $W$ 的右 $W_\Theta$-陪集，$\alpha\in\Pi$。则恰有三种可能：
\begin{enumerate}[label=(\roman*)]
\item $Cs_\alpha=C$；
\item $Cs_\alpha>C$，此时 $w_{Cs_\alpha}=w_Cs_\alpha$，且对任意 $w\in C$ 有 $\ell(ws_\alpha)=\ell(w)+1$；
\item $Cs_\alpha<C$，此时 $w_{Cs_\alpha}=w_Cs_\alpha$，且对任意 $w\in C$ 有 $\ell(ws_\alpha)=\ell(w)-1$。
\end{enumerate}
\end{proposition}

\begin{proof}
设 $Cs_\alpha\neq C$，$w$ 是 $C$ 的最短元。则有两种可能：$\ell(ws_\alpha)=\ell(w)+1$，或 $\ell(ws_\alpha)=\ell(w)-1$。

先设 $\ell(ws_\alpha)=\ell(w)+1$。取 $t\in W_\Theta$。若 $\ell(tws_\alpha)=\ell(tw)-1$，则由交换条件（注意 $\ell(ws_\alpha)=\ell(w)+1$）可知，存在 $t'\in W_\Theta$ 使 $ws_\alpha=t'w$。这会推出 $Cs_\alpha=C$，与假设矛盾。故 $\ell(tws_\alpha)=\ell(tw)+1$。这说明 $w_Cs_\alpha$ 是 $Cs_\alpha$ 的最长元，(ii) 成立。

现设 $\ell(ws_\alpha)=\ell(w)-1$。取 $t\in W_\Theta$。则
$$\ell(tws_\alpha)\le\ell(t)+\ell(ws_\alpha)=\ell(t)+\ell(w)-1=\ell(tw)-1,$$
故 $\ell(tws_\alpha)=\ell(tw)-1$。因此 $w_Cs_\alpha$ 是 $Cs_\alpha$ 的最长元，(iii) 成立。
\end{proof}

设 $C$ 是右 $W_\Theta$-陪集，$w$ 是它的最短元。则 $w_C=w_\Theta w$，且 $\ell(w_C)=\ell(w)+\ell(w_\Theta)$。这说明在 Bruhat 序下 $w_\Theta\le w_\Theta w=w_C$，即在 $W_\Theta\backslash W$ 的序中有 $W_\Theta\le C$。因此 $W_\Theta$ 是 $W_\Theta\backslash W$ 中的最小元。后面要用到下述刻画。

\begin{lemma}\label{lem:06.1.7}
设 $C\in W_\Theta\backslash W$。若对任意 $\alpha\in\Pi$，或者 $Cs_\alpha=C$，或者 $Cs_\alpha>C$，则 $C=W_\Theta$。
\end{lemma}

\begin{proof}
设 $w$ 是 $C$ 的最短元。假设说明，对任意 $\alpha\in\Pi$ 有 $\ell(ws_\alpha)=\ell(w)+1$。但这只可能在 $w=1$、$C=W_\Theta$ 时成立。
\end{proof}

\begin{lemma}\label{lem:06.1.8}
若 $w\in{}^{\Theta}W$，$t\in W_\Theta$，则
$$\ell(tw)=\ell(w)-\ell(t).$$
\end{lemma}

\begin{proof}
设 $w\in{}^{\Theta}W$，$t\in W_\Theta$。则 $w_\Theta w$ 是 $w$ 所在右 $W_\Theta$-陪集的最短元。又由 1.2，
\begin{align*}
\ell(tw)&=\ell((tw_\Theta)(w_\Theta w))=\ell(tw_\Theta)+\ell(w_\Theta w)\\
&=\ell(w_\Theta)-\ell(t)+\ell(w_\Theta w)=\ell(w)-\ell(t).
\end{align*}
\end{proof}

设 $B$ 是 $G$ 的 Borel 子群，$P_\Theta$ 是 $G$ 中含 $B$ 的型 $\Theta$ 的标准抛物子群。则 $X$ 中的 $P_\Theta$-轨道是 $B$-不变的，从而是 $X$ 中若干 Bruhat 胞腔的并。更确切地说，有下述结果。

\begin{lemma}\label{lem:06.1.9}
设 $O$ 是 $X$ 中的 $P_\Theta$-轨道，$C(w)\subset O$。则
$$O=\bigcup_{t\in W_\Theta}C(tw).$$
\end{lemma}

\begin{proof}
这由（[5], Ch.~IV, \S 2, no.~5, Prop.~2）得出。
\end{proof}

于是 $W_\Theta\backslash W$ 与 $X$ 中 $P_\Theta$-轨道的集之间有双射。设 $C$ 是 $W$ 的右 $W_\Theta$-陪集，$O$ 是相应的 $P_\Theta$-轨道。则由 1.8，
$$\dim O=\max_{t\in W_\Theta}\dim C(tw_C)=\max_{t\in W_\Theta}\ell(tw_C)=\ell(w_C).$$
故 $C(w_C)$ 是 $O$ 中的开 Bruhat 胞腔。这给出下述结果。

\begin{proposition}\label{prop:06.1.10}
对旗簇 $X$ 中的每个 $P_\Theta$-轨道 $O$，配上 $O$ 内唯一的开 Bruhat 胞腔 $C(w)$。这个对应是 $X$ 中全体 $P_\Theta$-轨道之集与全体 Bruhat 胞腔 $C(w)$（$w\in{}^{\Theta}W$）之集之间的双射。
\end{proposition}

最后给出 $W_\Theta\backslash W$ 上序关系的几何解释。

\begin{proposition}\label{prop:06.1.11}
设 $C\in W_\Theta\backslash W$，$O$ 是 $X$ 中对应于 $C$ 的 $P_\Theta$-轨道。则 $O$ 的闭包由 $X$ 中一切对应于 $D\le C$ 的 $P_\Theta$-轨道组成。
\end{proposition}

\begin{proof}
Bruhat 胞腔 $C(w_C)$ 在 $O$ 中开。因 $O$ 不可约，$C(w_C)$ 在 $O$ 中稠密，且 $\overline O=\overline{C(w_C)}$。故 $\overline O=\bigcup_{v\le w_C}C(v)$。另一方面，$\overline O$ 是若干 $P_\Theta$-轨道的并。设 $D$ 对应于 $\overline O$ 中的某个 $P_\Theta$-轨道。则 $C(w_D)$ 含于 $\overline O$，故 $w_D\le w_C$，即 $D\le C$。反之，若 $D\le C$，则 $w_D\le w_C$，且 $C(w_D)\subset\overline{C(w_C)}=\overline O$。若 $O'$ 是对应于 $D$ 的轨道，就得到
$$O'\subset\overline{O'}=\overline{C(w_D)}\subset\overline{C(w_C)}=\overline O.$$
\end{proof}

\section{广义 Verma 模}
\label{sec:06.6.2}

固定 $\g$ 的一个 Borel 子代数 $\bb_0$，并记 $B_0$ 为 $\operatorname{Int}(\g)$ 的单连通覆盖群 $G$ 中相应的 Borel 子群。对 $\Pi$ 的每个子集 $\Theta$，对应一个含 $\bb_0$ 的标准抛物子代数 $\p_\Theta$；记 $P_\Theta$ 为 $G$ 中相应的抛物子群。置 $\mathfrak{q}_\Theta=[\p_\Theta,\p_\Theta]$，并记 $Q_\Theta$ 为 $P_\Theta$ 的换位子群。

\begin{lemma}\label{lem:06.2.1}
\begin{enumerate}[label=(\roman*)]
\item $P_\Theta$ 的幂幺根 $N_\Theta$ 就是 $Q_\Theta$ 的幂幺根。
\item 设 $L_\Theta$ 是 $P_\Theta$ 的 Levi 因子。则 $L_\Theta$ 的换位子群 $S_\Theta$ 是 $Q_\Theta$ 的 Levi 因子。
\item $X$ 中的 $P_\Theta$-轨道也是 $Q_\Theta$-轨道。
\item 对 $x\in X$，$x$ 在 $P_\Theta$ 中的稳定化子连通。
\item 对 $x\in X$，$x$ 在 $Q_\Theta$ 中的稳定化子连通。
\end{enumerate}
\end{lemma}

\begin{proof}
设 $\mathfrak{c}$ 是 $\g$ 中含于 $\bb_0$ 的 Cartan 子代数，$R$ 是 $(\g,\mathfrak{c})$ 在 $\mathfrak{c}^*$ 中的根系。把 $\Theta$ 与 $R$ 中在 $\bb_0$ 决定的特殊化下对应于 $\Pi$ 的那组单根的一个子集等同起来。则 $\p_\Theta$ 由 $\mathfrak{c}$、正根对应的根子空间 $\g_\alpha$，以及形如 $-\beta$（$\beta$ 为 $\Theta$ 中若干根之和）的根对应的根子空间张成。

这说明 $N_\Theta$ 的 Lie 代数 $\n_\Theta$ 由那些「不是 $\Theta$ 中若干根之和」的正根对应的根子空间 $\g_\alpha$ 张成。Levi 因子 $L_\Theta$ 的 Lie 代数 $\lla_\Theta$ 则由 $\mathfrak{c}$ 与「$\Theta$ 中若干根之和及其相反数」对应的根子空间 $\g_\alpha$ 张成。置 $\s_\Theta=[\lla_\Theta,\lla_\Theta]$。因 $[\mathfrak{c},\n_\Theta]=\n_\Theta$，有
\begin{align*}
\mathfrak{q}_\Theta&=[\p_\Theta,\p_\Theta]=[\lla_\Theta,\lla_\Theta]+[\lla_\Theta,\n_\Theta]+[\n_\Theta,\n_\Theta]\\
&=\s_\Theta+\n_\Theta.
\end{align*}
由 Levi 分解的共轭性，这些结果与 $\mathfrak{c}$ 的选取无关。又 $P_\Theta$ 连通，故群层面的分解由 Lie 代数层面的结果立见。这就证完了 (i) 与 (ii)。

(iii) 易见 $Q_\Theta$ 与 $B_0$ 生成 $P_\Theta$。设 $O$ 是 $X$ 中的 $P_\Theta$-轨道，$x\in O$。则 $B_x\cap B_0$ 含 $G$ 的一个公共 Cartan 子群 $C$，而 $N_0$ 与 $C$ 生成 $B_0$。这说明 $Q_\Theta$ 与 $C$ 生成 $P_\Theta$，从而 $x$ 的 $Q_\Theta$-轨道与 $O$ 重合。

稳定化子 $P_\Theta\cap B_x$ 是 Cartan 子群 $C$ 与该稳定化子的幂幺根 $P_\Theta\cap N_x$ 的半直积。同样，稳定化子 $Q_\Theta\cap B_x$ 是 $S_\Theta$ 中某个 Cartan 子群与幂幺根 $Q_\Theta\cap N_x$ 的半直积。这立刻给出 (iv) 与 (v)。
\end{proof}

显然，对任意 $\Theta\subset\Pi$，$(\g,Q_\Theta)$ 是 Harish-Chandra 对。下面分析范畴 $\mathcal{M}_{coh}(\D_\lambda,Q_\Theta)$。先考虑标准 Harish-Chandra 层。

设 $\lambda\in\h^*$。固定 $X$ 中的一个 $P_\Theta$-轨道 $O$，$i_O:O\to X$ 是 $O$ 到 $X$ 的自然浸入。$X$ 上的齐性扭微分算子层 $\D_\lambda$ 在 $O$ 上定义了一个 $P_\Theta$-齐性扭微分算子层 $(\D_\lambda)_{i_O}$。由 2.1.(v)，$x$ 在 $Q_\Theta$ 中的稳定化子连通。因此 $O$ 上至多有一个不可约 $Q_\Theta$-齐性 $(\D_\lambda)_{i_O}$-联络。

置
$$P(\Sigma_\Theta)=\{\lambda\in\h^*\mid \check\alpha(\lambda)\in\Z,\ \alpha\in\Theta\}.$$

\begin{lemma}\label{lem:06.2.2}
设 $\lambda\in\h^*$。设 $O$ 是 $P_\Theta$-轨道，$w\in{}^{\Theta}W$ 使 Bruhat 胞腔 $C(w)$ 在 $O$ 中开。则下列条件等价：
\begin{enumerate}[label=(\roman*)]
\item $w\lambda\in P(\Sigma_\Theta)$；
\item $O$ 上存在不可约 $Q_\Theta$-齐性 $(\D_\lambda)_{i_O}$-联络。
\end{enumerate}
若这样的联络存在，则唯一。
\end{lemma}

\begin{proof}
设 $x\in C(w)$。则 $\bb_x$ 是 $\g$ 中相对于 $\bb_0$ 处于相对位置 $w$ 的 Borel 子代数。设 $\mathfrak{c}$ 是含于 $\bb_x\cap\bb_0$ 的 Cartan 子代数。关于 $\bb_0$ 决定的特殊化 $s$，$\n_x$ 由 $\alpha\in w(\Sigma^+)$ 对应的根子空间 $\g_\alpha$ 张成。由 2.1.(v)，$x$ 在 $Q_\Theta$ 中的稳定化子是以 $\mathfrak{q}_\Theta\cap\bb_x$ 为 Lie 代数的连通子群。因此关于这个特殊化，稳定化子的 Lie 代数由 $\mathfrak{q}_\Theta\cap\mathfrak{c}$ 与诸根子空间 $\g_\alpha$ 张成，其中 $\alpha$ 或者属于 $\Sigma^+\cap w(\Sigma^+)$，或者属于 $\Sigma_\Theta\cap w(\Sigma^+)$。由于 $w\in{}^{\Theta}W$，后一集合等于 $-\Sigma^+_\Theta$。

这说明 $x$ 的稳定化子是 $S_\Theta$ 中与 $B_0\cap S_\Theta$ 相对的 Borel 子群与正规子群 $N_0\cap N_x$ 的半直积。于是，$O$ 上一个不可约 $Q_\Theta$-齐性 $(\D_\lambda)_{i_O}$-联络由稳定化子的一个表示决定，而这个表示的微分必须来自 $P(\Sigma_\Theta)$ 中某个元素的特殊化在稳定化子上的限制。由 $\bb_x$ 定义的相对化 $s'$ 等于 $s\circ w$。故决定该联络的线性形式是 $s'(\lambda+\rho)=s(w\lambda+w\rho)$，这说明 $w\lambda$ 必须属于 $P(\Sigma_\Theta)$。
\end{proof}

这说明，对任意 $\lambda\in\h^*$，$\mathcal{M}_{coh}(\D_\lambda,Q_\Theta)$ 中附着于轨道 $O$ 的标准 Harish-Chandra 层至多有一个。记之为 $I(O,\lambda)$，其唯一的不可约 Harish-Chandra 子层记为 $L(O,\lambda)$。先注意，不可约 Harish-Chandra 层 $L(O,\lambda)$ 其实同构于前面遇到的不可约模。更确切地说，有下述结果。

\begin{proposition}\label{prop:06.2.3}
设 $O$ 是 $X$ 中的 $P_\Theta$-轨道，$C(w)$ 是 $O$ 中开的 Bruhat 胞腔。设 $w\lambda\in P(\Sigma_\Theta)$。则 $L(O,\lambda)=L(w,\lambda)$。
\end{proposition}

\begin{proof}
设 $j:C(w)\to O$ 是自然浸入，$\tau$ 是 $O$ 上唯一的不可约 $Q_\Theta$-齐性 $(\D_\lambda)_{i_O}$-联络。则 $\tau$ 在 $C(w)$ 上的限制是不可约 $N_0$-联络，从而同构于 $\calO_{C(w)}$。这说明 $\tau\subset R^0j_+(\calO_{C(w)})$。于是
\begin{align*}
I(O,\lambda)&=R^0i_{O+}(\tau)\subset R^0i_{O+}(R^0j_+(\calO_{C(w)}))\\
&=R^0i_{w+}(\calO_{C(w)})=I(w,\lambda).
\end{align*}
故 $L(O,\lambda)\subset I(w,\lambda)$。由于 $L(O,\lambda)$ 不可约，它必等于 $L(w,\lambda)$。
\end{proof}

下面分析 $I(O,\lambda)$。我们沿用 V.2 的方法。
